% generated by GAPDoc2LaTeX from XML source (Frank Luebeck)
\documentclass[a4paper,11pt]{report}

\usepackage[top=37mm,bottom=37mm,left=27mm,right=27mm]{geometry}
\sloppy
\pagestyle{myheadings}
\usepackage{amssymb}
\usepackage[utf8]{inputenc}
\usepackage{makeidx}
\makeindex
\usepackage{color}
\definecolor{FireBrick}{rgb}{0.5812,0.0074,0.0083}
\definecolor{RoyalBlue}{rgb}{0.0236,0.0894,0.6179}
\definecolor{RoyalGreen}{rgb}{0.0236,0.6179,0.0894}
\definecolor{RoyalRed}{rgb}{0.6179,0.0236,0.0894}
\definecolor{LightBlue}{rgb}{0.8544,0.9511,1.0000}
\definecolor{Black}{rgb}{0.0,0.0,0.0}

\definecolor{linkColor}{rgb}{0.0,0.0,0.554}
\definecolor{citeColor}{rgb}{0.0,0.0,0.554}
\definecolor{fileColor}{rgb}{0.0,0.0,0.554}
\definecolor{urlColor}{rgb}{0.0,0.0,0.554}
\definecolor{promptColor}{rgb}{0.0,0.0,0.589}
\definecolor{brkpromptColor}{rgb}{0.589,0.0,0.0}
\definecolor{gapinputColor}{rgb}{0.589,0.0,0.0}
\definecolor{gapoutputColor}{rgb}{0.0,0.0,0.0}

%%  for a long time these were red and blue by default,
%%  now black, but keep variables to overwrite
\definecolor{FuncColor}{rgb}{0.0,0.0,0.0}
%% strange name because of pdflatex bug:
\definecolor{Chapter }{rgb}{0.0,0.0,0.0}
\definecolor{DarkOlive}{rgb}{0.1047,0.2412,0.0064}


\usepackage{fancyvrb}

\usepackage{mathptmx,helvet}
\usepackage[T1]{fontenc}
\usepackage{textcomp}


\usepackage[
            pdftex=true,
            bookmarks=true,        
            a4paper=true,
            pdftitle={Written with GAPDoc},
            pdfcreator={LaTeX with hyperref package / GAPDoc},
            colorlinks=true,
            backref=page,
            breaklinks=true,
            linkcolor=linkColor,
            citecolor=citeColor,
            filecolor=fileColor,
            urlcolor=urlColor,
            pdfpagemode={UseNone}, 
           ]{hyperref}

\newcommand{\maintitlesize}{\fontsize{50}{55}\selectfont}

% write page numbers to a .pnr log file for online help
\newwrite\pagenrlog
\immediate\openout\pagenrlog =\jobname.pnr
\immediate\write\pagenrlog{PAGENRS := [}
\newcommand{\logpage}[1]{\protect\write\pagenrlog{#1, \thepage,}}
%% were never documented, give conflicts with some additional packages

\newcommand{\GAP}{\textsf{GAP}}

%% nicer description environments, allows long labels
\usepackage{enumitem}
\setdescription{style=nextline}

%% depth of toc
\setcounter{tocdepth}{1}





%% command for ColorPrompt style examples
\newcommand{\gapprompt}[1]{\color{promptColor}{\bfseries #1}}
\newcommand{\gapbrkprompt}[1]{\color{brkpromptColor}{\bfseries #1}}
\newcommand{\gapinput}[1]{\color{gapinputColor}{#1}}


\begin{document}

\logpage{[ 0, 0, 0 ]}
\begin{titlepage}
\mbox{}\vfill

\begin{center}{\maintitlesize \textbf{ hyparr \mbox{}}}\\
\vfill

\hypersetup{pdftitle= hyparr }
\markright{\scriptsize \mbox{}\hfill  hyparr  \hfill\mbox{}}
{\Huge \textbf{ Computations with hyperplane arrangements. \mbox{}}}\\
\vfill

{\Huge  0.1 \mbox{}}\\[1cm]
{ 13 March 2026 \mbox{}}\\[1cm]
\mbox{}\\[2cm]
{\Large \textbf{ Paul Muecksch\\
    \mbox{}}}\\
\hypersetup{pdfauthor= Paul Muecksch\\
    }
\end{center}\vfill

\mbox{}\\
{\mbox{}\\
\small \noindent \textbf{ Paul Muecksch\\
    }  Email: \href{mailto://paul.muecksch+uni@gmail.com} {\texttt{paul.muecksch+uni@gmail.com}}\\
  Homepage: \href{https://paulmuecksch.github.io} {\texttt{https://paulmuecksch.github.io}}\\
  Address: \begin{minipage}[t]{8cm}\noindent
 Welfengarten 1, 30167 Hannover, Germany\\
 \end{minipage}
}\\
\end{titlepage}

\newpage\setcounter{page}{2}
\newpage

\def\contentsname{Contents\logpage{[ 0, 0, 1 ]}}

\tableofcontents
\newpage

     
\chapter{\textcolor{Chapter }{Introduction}}\label{Chapter_Introduction}
\logpage{[ 1, 0, 0 ]}
\hyperdef{L}{X7DFB63A97E67C0A1}{}
{
  

 HypArr is a package for computations with hyperplane arrangements, oriented
matroids, and their topological invariants, in particular Milnor fibers and
complements. 

 HypArr is a package for computations with hyperplane arrangements, oriented
matroids, and their topological invariants, in particular Milnor fibers and
complements. 

 }

   
\chapter{\textcolor{Chapter }{Hyperplane arrangements}}\label{Chapter_Hyperplane_arrangements}
\logpage{[ 2, 0, 0 ]}
\hyperdef{L}{X7EF7AC44814961EC}{}
{
  
\section{\textcolor{Chapter }{Attributes}}\label{Chapter_Hyperplane_arrangements_Section_Attributes}
\logpage{[ 2, 1, 0 ]}
\hyperdef{L}{X7C701DBF7BAE649A}{}
{
  

\subsection{\textcolor{Chapter }{Dimension (for IsHyperplaneArrangement)}}
\logpage{[ 2, 1, 1 ]}\nobreak
\hyperdef{L}{X7CCC1A2A7C3F2CC5}{}
{\noindent\textcolor{FuncColor}{$\triangleright$\enspace\texttt{Dimension({\mdseries\slshape A})\index{Dimension@\texttt{Dimension}!for IsHyperplaneArrangement}
\label{Dimension:for IsHyperplaneArrangement}
}\hfill{\scriptsize (attribute)}}\\
\textbf{\indent Returns:\ }
a non\texttt{\symbol{45}}negative integer. 



 Returns the dimension of the ambient vector space in which the hyperplane
arrangement $\mathcal{A}$ is defined. }

 

\subsection{\textcolor{Chapter }{Roots (for IsHyperplaneArrangement)}}
\logpage{[ 2, 1, 2 ]}\nobreak
\hyperdef{L}{X8419FA5F7A08DF56}{}
{\noindent\textcolor{FuncColor}{$\triangleright$\enspace\texttt{Roots({\mdseries\slshape A})\index{Roots@\texttt{Roots}!for IsHyperplaneArrangement}
\label{Roots:for IsHyperplaneArrangement}
}\hfill{\scriptsize (attribute)}}\\
\textbf{\indent Returns:\ }
A list of vectors defining the hyperplanes of \mbox{\texttt{\mdseries\slshape A}}. 



 Returns the list of defining linear forms (also called \emph{roots}) of the hyperplanes in the arrangement \mbox{\texttt{\mdseries\slshape A}}. Each vector represents a hyperplane by the linear equation 
\[ r_1 x_1 + \cdots + r_n x_n = 0. \]
 }

 

\subsection{\textcolor{Chapter }{IntersectionLattice (for IsHyperplaneArrangement)}}
\logpage{[ 2, 1, 3 ]}\nobreak
\hyperdef{L}{X8555074B83058460}{}
{\noindent\textcolor{FuncColor}{$\triangleright$\enspace\texttt{IntersectionLattice({\mdseries\slshape A})\index{IntersectionLattice@\texttt{IntersectionLattice}!for IsHyperplaneArrangement}
\label{IntersectionLattice:for IsHyperplaneArrangement}
}\hfill{\scriptsize (attribute)}}\\
\textbf{\indent Returns:\ }
A list of lists describing the intersections of \mbox{\texttt{\mdseries\slshape A}}. 



 Computes the intersection lattice of the hyperplane arrangement \mbox{\texttt{\mdseries\slshape A}}. 

 The lattice is represented
level\texttt{\symbol{45}}by\texttt{\symbol{45}}level. The entry $L[k]$ contains all intersections of $k-1$ hyperplanes. 

 Each lattice element is stored as a list of indices referring to hyperplanes
in \texttt{Roots(A)}. 

 For example, the list \texttt{[1,3]} represents the intersection of the first and third hyperplane of the
arrangement. 

 The algorithm incrementally constructs the lattice by adding hyperplanes one
by one. 

 
\begin{Verbatim}[commandchars=!@|,fontsize=\small,frame=single,label=Example]
  !gapprompt@gap>| !gapinput@A := HyperplaneArrangement([[1,0],[0,1],[1,1]]);|
   gap> IntersectionLattice(A);
   [ [ [ 1 ], [ 2 ], [ 3 ] ],
     [ [ 1, 2 ], [ 1, 3 ], [ 2, 3 ] ],
     [ [ 1, 2, 3 ] ] ]
\end{Verbatim}
 }

 

\subsection{\textcolor{Chapter }{MSetInvL (for IsHyperplaneArrangement)}}
\logpage{[ 2, 1, 4 ]}\nobreak
\hyperdef{L}{X852CC1A8795FCD15}{}
{\noindent\textcolor{FuncColor}{$\triangleright$\enspace\texttt{MSetInvL({\mdseries\slshape A})\index{MSetInvL@\texttt{MSetInvL}!for IsHyperplaneArrangement}
\label{MSetInvL:for IsHyperplaneArrangement}
}\hfill{\scriptsize (attribute)}}\\
\textbf{\indent Returns:\ }
A list encoding multiset invariants of the intersection lattice. 



 Computes multiset invariants of the intersection lattice of the hyperplane
arrangement \mbox{\texttt{\mdseries\slshape A}}. 

 For each level of the lattice, the function records how many intersections
occur with a given number of defining hyperplanes. 

 These invariants can be used to compare arrangements or detect combinatorial
similarities between intersection lattices. }

 

\subsection{\textcolor{Chapter }{AGpql}}
\logpage{[ 2, 1, 5 ]}\nobreak
\hyperdef{L}{X858205A37B7EB5FA}{}
{\noindent\textcolor{FuncColor}{$\triangleright$\enspace\texttt{AGpql({\mdseries\slshape Intgers, p, q, l})\index{AGpql@\texttt{AGpql}}
\label{AGpql}
}\hfill{\scriptsize (function)}}\\
\textbf{\indent Returns:\ }
A hyperplane arrangement. 



 Constructs the reflection arrangement associated to the monomial complex
reflection group $G(p,q,l)$. The hyperplanes are defined by equations of the form 
\[ x_i = \zeta^k x_j \]
 where $\zeta$ is a primitive $p$\texttt{\symbol{45}}th root of unity. 
\begin{Verbatim}[commandchars=!@|,fontsize=\small,frame=single,label=Example]
  !gapprompt@gap>| !gapinput@A := AGpql(2,1,3);|
   <HyperplaneArrangement: 9 hyperplanes in 3-space>
  !gapprompt@gap>| !gapinput@Roots(A);|
  [ [ 1, 0, 0 ], [ 0, 1, 0 ], [ 0, 0, 1 ], [ 1, 1, 0 ], [ 1, -1, 0 ], [ 1, 0, 1 ], [ 1, 0, -1 ], [ 0, 1, 1 ], [ 0, 1, -1 ] ]
  
\end{Verbatim}
 }

 

\subsection{\textcolor{Chapter }{HArrResHind}}
\logpage{[ 2, 1, 6 ]}\nobreak
\hyperdef{L}{X7CC338E67CE223FE}{}
{\noindent\textcolor{FuncColor}{$\triangleright$\enspace\texttt{HArrResHind({\mdseries\slshape A})\index{HArrResHind@\texttt{HArrResHind}}
\label{HArrResHind}
}\hfill{\scriptsize (function)}}\\
\textbf{\indent Returns:\ }
A hyperplane arrangement. 



 Computes the restriction of the arrangement \mbox{\texttt{\mdseries\slshape A}} to the \mbox{\texttt{\mdseries\slshape k}}\texttt{\symbol{45}}th hyperplane of \texttt{Roots(A)}. }

 

\subsection{\textcolor{Chapter }{HArrResX}}
\logpage{[ 2, 1, 7 ]}\nobreak
\hyperdef{L}{X784AA86286AA9DD8}{}
{\noindent\textcolor{FuncColor}{$\triangleright$\enspace\texttt{HArrResX({\mdseries\slshape arg})\index{HArrResX@\texttt{HArrResX}}
\label{HArrResX}
}\hfill{\scriptsize (function)}}\\


 

 }

 

\subsection{\textcolor{Chapter }{Essentialization}}
\logpage{[ 2, 1, 8 ]}\nobreak
\hyperdef{L}{X8221BCFC7DE0069A}{}
{\noindent\textcolor{FuncColor}{$\triangleright$\enspace\texttt{Essentialization({\mdseries\slshape A})\index{Essentialization@\texttt{Essentialization}}
\label{Essentialization}
}\hfill{\scriptsize (function)}}\\
\textbf{\indent Returns:\ }
A hyperplane arrangement. 



 Returns the \emph{essentialization} of the hyperplane arrangement An arrangement is called essential if the
intersection of all its hyperplanes is the origin. If this is not the case,
the function restricts the arrangement to a complementary subspace so that the
resulting arrangement becomes essential. }

 }

 
\section{\textcolor{Chapter }{Constructors}}\label{Chapter_Hyperplane_arrangements_Section_Constructors}
\logpage{[ 2, 2, 0 ]}
\hyperdef{L}{X86EC0F0A78ECBC10}{}
{
  

\subsection{\textcolor{Chapter }{HyperplaneArrangement (for IsList)}}
\logpage{[ 2, 2, 1 ]}\nobreak
\hyperdef{L}{X860548D886CE5BA2}{}
{\noindent\textcolor{FuncColor}{$\triangleright$\enspace\texttt{HyperplaneArrangement({\mdseries\slshape r})\index{HyperplaneArrangement@\texttt{HyperplaneArrangement}!for IsList}
\label{HyperplaneArrangement:for IsList}
}\hfill{\scriptsize (operation)}}\\
\textbf{\indent Returns:\ }
A hyperplane arrangement 



 Constructs a hyperplane arrangement from a list \mbox{\texttt{\mdseries\slshape R}} of vectors representing defining linear forms of hyperplanes. 

 Each vector $r = [r_1,\dots,r_n]$ corresponds to the hyperplane 
\[ r_1 x_1 + \cdots + r_n x_n = 0. \]
 The vectors must lie in the same vector space. 

 Linearly dependent defining forms are removed automatically, so the resulting
arrangement only stores pairwise linearly independent hyperplanes. 

 
\begin{Verbatim}[commandchars=!@|,fontsize=\small,frame=single,label=Example]
  !gapprompt@gap>| !gapinput@A := HyperplaneArrangement([[1,0,0],[0,1,0],[0,0,1]]);|
  <HyperplaneArrangement: 3 hyperplanes in 3-space>
  !gapprompt@gap>| !gapinput@Roots(A);|
  [ [ 1, 0, 0 ], [ 0, 1, 0 ], [ 0, 0, 1 ] ]
\end{Verbatim}
 }

 }

 }

   
\chapter{\textcolor{Chapter }{Oriented Matroids}}\label{Chapter_Oriented_Matroids}
\logpage{[ 3, 0, 0 ]}
\hyperdef{L}{X8388FBA579602CDB}{}
{
  

 HypArr is a package for computations with hyperplane arrangements, oriented
matroids, and their topological invariants, in particular Milnor fibers and
complements. 

 HypArr is a package for computations with hyperplane arrangements, oriented
matroids, and their topological invariants, in particular Milnor fibers and
complements. 

 }

   
\chapter{\textcolor{Chapter }{Functionality}}\label{Chapter_Functionality}
\logpage{[ 4, 0, 0 ]}
\hyperdef{L}{X87F1120883F5B4D0}{}
{
  

 
\section{\textcolor{Chapter }{Example Methods}}\label{Chapter_Functionality_Section_Example_Methods}
\logpage{[ 4, 1, 0 ]}
\hyperdef{L}{X7ED1D62579C10890}{}
{
  

 This section will describe the example methods of HypArr 

 This section will describe the example methods of HypArr 

\subsection{\textcolor{Chapter }{HypARr{\textunderscore}Example}}
\logpage{[ 4, 1, 1 ]}\nobreak
\hyperdef{L}{X7F57601A8251E180}{}
{\noindent\textcolor{FuncColor}{$\triangleright$\enspace\texttt{HypARr{\textunderscore}Example({\mdseries\slshape arg})\index{HypARr{\textunderscore}Example@\texttt{HypARr{\textunderscore}Example}}
\label{HypARruScoreExample}
}\hfill{\scriptsize (function)}}\\


 Insert documentation for your function here }

 

 This section will describe the example methods of HypArr }

 

 }

 \def\indexname{Index\logpage{[ "Ind", 0, 0 ]}
\hyperdef{L}{X83A0356F839C696F}{}
}

\cleardoublepage
\phantomsection
\addcontentsline{toc}{chapter}{Index}


\printindex

\immediate\write\pagenrlog{["Ind", 0, 0], \arabic{page},}
\newpage
\immediate\write\pagenrlog{["End"], \arabic{page}];}
\immediate\closeout\pagenrlog
\end{document}
